% TOP_PROBLEM: {"id": 1197, "title": "The Sofic Groups Conjecture", "category_id": 4, "category": {"id": 4, "name": "algebra", "display_name": "Algebra", "description": "Group theory, ring theory, field theory, and algebraic structures.", "slug": "algebra", "order_index": 4, "created_at": "2026-07-31T15:26:25.671Z"}, "status": "open"}
% FIRST_POSTED: 2026-10-06
% ATTEMPT_STATUS: solved
% OPENAI_CLAIM: {"schema_version":1,"source_repo":"https://github.com/openai/math","source_commit":"adc7f1241b42e322a6451854ab7e4b4c146bf78a","release_date":"2026-10-06","match":"full","resolution":"disproved","independently_reviewed":false,"second_pass":"agreed","families":[{"family":"197","title":"A torsion-free group algebra that is not directly finite","subject":"Algebra","manuscripts":[{"title":"A Counterexample to Kaplansky's Direct-Finiteness Conjecture in Characteristic Two","pdf_path":"preprints/A-Counterexample-to-Kaplanskys-Direct-Finiteness-Conjecture-in-Characteristic-Two-September-23-2026/paper.pdf","readme_path":"preprints/A-Counterexample-to-Kaplanskys-Direct-Finiteness-Conjecture-in-Characteristic-Two-September-23-2026/README.md","date":"2026-09-23","theorem_ref":"thm:main (sections/01-introduction.tex:8)","statement_sources":[{"path":"preprints/A-Counterexample-to-Kaplanskys-Direct-Finiteness-Conjecture-in-Characteristic-Two-September-23-2026/build/sections/01-introduction.tex","line":8,"label":"thm:main"}]},{"title":"A Counterexample to Kaplansky's Direct-Finiteness Conjecture in Odd Characteristic","pdf_path":"preprints/A-Counterexample-to-Kaplanskys-Direct-Finiteness-Conjecture-in-Odd-Characteristic-September-26-2026/paper.pdf","readme_path":"preprints/A-Counterexample-to-Kaplanskys-Direct-Finiteness-Conjecture-in-Odd-Characteristic-September-26-2026/README.md","date":"2026-09-26","theorem_ref":"op:main (sections/01-introduction.tex:13)","statement_sources":[{"path":"preprints/A-Counterexample-to-Kaplanskys-Direct-Finiteness-Conjecture-in-Odd-Characteristic-September-26-2026/build/sections/01-introduction.tex","line":13,"label":"op:main"}]},{"title":"A Counterexample to the Group-Ring Determinant Conjecture","pdf_path":"preprints/A-Counterexample-to-the-Group-Ring-Determinant-Conjecture-September-23-2026/paper.pdf","readme_path":"preprints/A-Counterexample-to-the-Group-Ring-Determinant-Conjecture-September-23-2026/README.md","date":"2026-09-23","theorem_ref":"thm:main (main.tex:90); thm:input (main.tex:107)","statement_sources":[{"path":"preprints/A-Counterexample-to-the-Group-Ring-Determinant-Conjecture-September-23-2026/build/main.tex","line":90,"label":"thm:main"},{"path":"preprints/A-Counterexample-to-the-Group-Ring-Determinant-Conjecture-September-23-2026/build/main.tex","line":107,"label":"thm:input"}]},{"title":"A Torsion-Free Group Algebra That Is Not Directly Finite","pdf_path":"preprints/A-Torsion-Free-Group-Algebra-That-Is-Not-Directly-Finite-October-4-2026/direct-finiteness.pdf","readme_path":"preprints/A-Torsion-Free-Group-Algebra-That-Is-Not-Directly-Finite-October-4-2026/README.md","date":"2026-10-04","theorem_ref":"thm:main (sections/introduction.tex:13)","statement_sources":[{"path":"preprints/A-Torsion-Free-Group-Algebra-That-Is-Not-Directly-Finite-October-4-2026/build/sections/introduction.tex","line":13,"label":"thm:main"}]}],"lean":{"doc_path":"lean/docs/197.md","comparators":["lean/ComparatorChallenges/GroupRingDeterminant.json","lean/ComparatorChallenges/GroupRingDeterminant.lean","lean/ComparatorChallenges/KaplanskyDirectFiniteness.json","lean/ComparatorChallenges/KaplanskyDirectFiniteness.lean","lean/ComparatorChallenges/KaplanskyFinitelyPresented.json","lean/ComparatorChallenges/KaplanskyFinitelyPresented.lean","lean/ComparatorChallenges/OddKaplansky.json","lean/ComparatorChallenges/OddKaplansky.lean"],"declarations":["OAI.GroupRingDeterminant.main","OAI.KaplanskyCounterexample.finitelyPresented_counterexample","OAI.OddKaplansky.main_theorem"],"covers_main_theorem":false},"reasoning_trace":"reasoning_traces/kaplansky-direct-finiteness-characteristic-two.pdf"}]}
\documentclass[11pt]{article}
\usepackage[margin=25mm]{geometry}
\usepackage{fontspec}
\setmainfont{Latin Modern Roman}
\usepackage{amsmath,amssymb,mathtools}
\usepackage{unicode-math}
\setmathfont{Latin Modern Math}
\usepackage{enumitem}
\usepackage{xurl,hyperref,bookmark}
\hypersetup{colorlinks=true,urlcolor=blue}
\setcounter{secnumdepth}{0}
\setlength{\parindent}{0pt}
\setlength{\parskip}{5pt}
\setlength{\emergencystretch}{3em}
\begin{document}
\section*{The Sofic Groups Conjecture}

\subsection{OpenAI's claim}
The linked manuscripts present the following claims by OpenAI: A finite characteristic-two field and a finitely presented group are claimed to produce a scalar one-sided inverse pair; the resulting shift is injective without being onto and the group is nonsofic; For a particular odd prime $p$, the authors claim a one-sided inverse defect over a field with $p^4$ elements, together with a nonsurjective injective shift; An integral group-ring matrix is claimed to be rationally invertible while its Fuglede--Kadison determinant lies strictly between $0$ and $1$; Another construction claims a torsion-free, finitely presented $G$ for which $\mathbb F_2[G]$ contains $a,b$ satisfying $ab=1$ and $ba\ne1$. Together they give a claimed negative answer in the matching scope.

\subsection{Mathematical statement of the claimed result}
\paragraph{Result 1.} Take a finite field \(K\) with \(\operatorname{char}K=2\). The constructed group \(G\) is finitely presented and has an element of odd prime order. Two finite group-ring sums meet
\[
 a_{\mathrm{out}},b_{\mathrm{out}}\in K[G],\qquad
 a_{\mathrm{out}}b_{\mathrm{out}}=1,\qquad
 b_{\mathrm{out}}a_{\mathrm{out}}\ne1.
\]
Its specification uses terminating finite-set and finite-field computations. Unlike the torsion-free companion, this particular group contains torsion.

\paragraph{Result 2.} Define the integers
\[
 m=\binom{1200}{600},\qquad
 p=\min\{h\in\mathbb Z:h>1,\ h\mid(m!)^2+1\}.
\]
The claimed example has \(p\) an odd prime, \(|K|=p^4\), and a specified finitely generated group \(G\) containing torsion. Its finite sums \(a,b\in K[G]\) satisfy
\[
 ab=1\quad\hbox{and}\quad ba\ne1.
\]
This is an example in a specified odd characteristic, rather than a theorem for each odd prime.

\paragraph{Result 3.} For some finitely generated discrete group \(G_D\), integer \(n\ge1\), and matrix
\[
 A\in M_n(\mathbb Z[G_D]),\qquad A\text{ invertible over }\mathbb Q[G_D],
\]
the operator induced by the left regular action \(T_A\) on \(\ell^2(G_D)^n\) has Fuglede--Kadison determinant
\[
 0<\det_{\mathcal N(G_D)}(T_A)<1.
\]
Here \(\mathcal N(G_D)\) is the group von Neumann algebra and the logarithmic spectral integral defining the determinant is finite and negative. The characteristic-two direct-finiteness result quoted in the proof is an input, not an additional determinant hypothesis. The convention is $(L_g\xi)(h)=\xi(g^{-1}h)$, and the matrix trace $\tau_n$ is unnormalized. For the normalized spectral law $\mu_A$ of $T_A^*T_A$,
\[
 \log\det_{\mathcal N(G_D)}(T_A)=\tfrac12\tau_n(\log(T_A^*T_A))=\tfrac n2\int\log t\,d\mu_A(t).
\]

\paragraph{Result 4.} The construction supplies a torsion-free, finitely presented group \(G\), with a finite two-dimensional \(K(G,1)\), and elements of its characteristic-two group algebra satisfying
\[
 a,b,c\in\mathbb F_2[G],\qquad ab=1,\quad ac=0,\quad c\ne0,\quad ba\ne1.
\]
In particular, \(a\ne0\) and \(e=ba\) obeys \(e^2=e\), \(e\notin\{0,1\}\). Thus this same example supplies failures of direct finiteness, absence of zero divisors, and triviality of scalar idempotents. These assertions concern \(\mathbb F_2[G]\); the separate integral-unit assertion is not a consequence.

\subsection{Outline of the argument}
Incidence data and polynomial pairings first give a defective inverse pair of matrices. A geometric group-embedding construction preserves the defect, and character packing realizes the matrix algebra inside a larger group algebra; dual convolution gives the cellular-automaton consequence. The construction uses translation representations of finite groups: a high-degree subspace is killed on one side and a weighted pairing supplies the other annihilation. Finite group identifications and a quadratic field tower package these local relations into the scalar inverse pair; the claim concerns the specified prime rather than all odd characteristics.

A characteristic-two inverse defect is lifted using finite-coordinate projections in a wreath-product setting. These projections detect a selected coefficient, and a tracial logarithmic derivative applied to block shears calculates a determinant below the conjectured integral lower bound. An unequal-side parity pattern supplies the one-sided inverse relation in labeled graphs. Weighted counts handle repeated edge traversals, and a separator reduces arbitrarily large planar arrangements to bounded path collections; cone attachment and root separation then preserve the inverse defect while proving torsion-freeness.

\subsection{Scope relative to this problem}
The main finite-field one-sided-inverse pair produces a nonsurjunctive group. Soficity entails surjunctivity and direct finiteness for the relevant finite-field algebra, so this same finitely generated group is not sofic. For the integral matrix of Theorem 1.1, the Fuglede--Kadison determinant belongs to $(0,1)$. Sofic groups satisfy the determinant conjecture for every integral group-ring matrix, so the example’s group cannot be sofic. Within these assumptions, this is a claimed negative answer to the stated question.

\subsection{Formal verification}
For family 197, the documented Lean entries include the scalar inverse defects in the specified characteristics and an operator-invertible determinant counterexample. The characteristic-two search recipe is included, but the nonsofic consequence, the extra spectral-integral conclusion, and the later torsion-free construction are outside the selected statements. The listed Comparator files are \texttt{GroupRingDeterminant.lean}, \texttt{KaplanskyDirectFiniteness.lean}, \texttt{KaplanskyFinitelyPresented.lean}, \texttt{OddKaplansky.lean}. Named declarations include \texttt{OAI.GroupRingDeterminant.main}, \texttt{OAI.KaplanskyCounterexample.finitelyPresented\_counterexample}, \texttt{OAI.OddKaplansky.main\_theorem}. This site has not executed the Lean code. The scope descriptions do not certify the remaining manuscript arguments.

\subsection{Status caveat}
OpenAI attributes these manuscripts to work by an internal model. Its release cautions that arguments without formalization can contain mistakes. This catalogue has not independently checked the proofs; the labels report the mathematical claims attributed to OpenAI.

\subsection{Sources}
\begin{itemize}
\item OpenAI, \emph{A Counterexample to Kaplansky's Direct-Finiteness Conjecture in Characteristic Two}, 2026-09-23; citation key \texttt{OAI:A-Counterexample-to-Kaplanskys-Direct-Finiteness-Conjecture-in-Characteristic-Two-September-23-2026}. \href{https://github.com/openai/math/blob/adc7f1241b42e322a6451854ab7e4b4c146bf78a/preprints/A-Counterexample-to-Kaplanskys-Direct-Finiteness-Conjecture-in-Characteristic-Two-September-23-2026/paper.pdf}{PDF (GitHub)}; \href{https://github.com/openai/math/raw/adc7f1241b42e322a6451854ab7e4b4c146bf78a/preprints/A-Counterexample-to-Kaplanskys-Direct-Finiteness-Conjecture-in-Characteristic-Two-September-23-2026/paper.pdf}{PDF (direct)}; \href{https://github.com/openai/math/blob/adc7f1241b42e322a6451854ab7e4b4c146bf78a/preprints/A-Counterexample-to-Kaplanskys-Direct-Finiteness-Conjecture-in-Characteristic-Two-September-23-2026/README.md}{Manuscript page and citation}; \href{https://github.com/openai/math/blob/main/preprints/A-Counterexample-to-Kaplanskys-Direct-Finiteness-Conjecture-in-Characteristic-Two-September-23-2026/paper.pdf}{Latest version}; Statement source: \href{https://github.com/openai/math/blob/adc7f1241b42e322a6451854ab7e4b4c146bf78a/preprints/A-Counterexample-to-Kaplanskys-Direct-Finiteness-Conjecture-in-Characteristic-Two-September-23-2026/build/sections/01-introduction.tex\#L8}{thm:main}.
\item OpenAI, \emph{A Counterexample to Kaplansky's Direct-Finiteness Conjecture in Odd Characteristic}, 2026-09-26; citation key \texttt{OAI:A-Counterexample-to-Kaplanskys-Direct-Finiteness-Conjecture-in-Odd-Characteristic-September-26-2026}. \href{https://github.com/openai/math/blob/adc7f1241b42e322a6451854ab7e4b4c146bf78a/preprints/A-Counterexample-to-Kaplanskys-Direct-Finiteness-Conjecture-in-Odd-Characteristic-September-26-2026/paper.pdf}{PDF (GitHub)}; \href{https://github.com/openai/math/raw/adc7f1241b42e322a6451854ab7e4b4c146bf78a/preprints/A-Counterexample-to-Kaplanskys-Direct-Finiteness-Conjecture-in-Odd-Characteristic-September-26-2026/paper.pdf}{PDF (direct)}; \href{https://github.com/openai/math/blob/adc7f1241b42e322a6451854ab7e4b4c146bf78a/preprints/A-Counterexample-to-Kaplanskys-Direct-Finiteness-Conjecture-in-Odd-Characteristic-September-26-2026/README.md}{Manuscript page and citation}; \href{https://github.com/openai/math/blob/main/preprints/A-Counterexample-to-Kaplanskys-Direct-Finiteness-Conjecture-in-Odd-Characteristic-September-26-2026/paper.pdf}{Latest version}; Statement source: \href{https://github.com/openai/math/blob/adc7f1241b42e322a6451854ab7e4b4c146bf78a/preprints/A-Counterexample-to-Kaplanskys-Direct-Finiteness-Conjecture-in-Odd-Characteristic-September-26-2026/build/sections/01-introduction.tex\#L13}{op:main}.
\item OpenAI, \emph{A Counterexample to the Group-Ring Determinant Conjecture}, 2026-09-23; citation key \texttt{OAI:A-Counterexample-to-the-Group-Ring-Determinant-Conjecture-September-23-2026}. \href{https://github.com/openai/math/blob/adc7f1241b42e322a6451854ab7e4b4c146bf78a/preprints/A-Counterexample-to-the-Group-Ring-Determinant-Conjecture-September-23-2026/paper.pdf}{PDF (GitHub)}; \href{https://github.com/openai/math/raw/adc7f1241b42e322a6451854ab7e4b4c146bf78a/preprints/A-Counterexample-to-the-Group-Ring-Determinant-Conjecture-September-23-2026/paper.pdf}{PDF (direct)}; \href{https://github.com/openai/math/blob/adc7f1241b42e322a6451854ab7e4b4c146bf78a/preprints/A-Counterexample-to-the-Group-Ring-Determinant-Conjecture-September-23-2026/README.md}{Manuscript page and citation}; \href{https://github.com/openai/math/blob/main/preprints/A-Counterexample-to-the-Group-Ring-Determinant-Conjecture-September-23-2026/paper.pdf}{Latest version}; Statement source: \href{https://github.com/openai/math/blob/adc7f1241b42e322a6451854ab7e4b4c146bf78a/preprints/A-Counterexample-to-the-Group-Ring-Determinant-Conjecture-September-23-2026/build/main.tex\#L90}{thm:main}; Statement source: \href{https://github.com/openai/math/blob/adc7f1241b42e322a6451854ab7e4b4c146bf78a/preprints/A-Counterexample-to-the-Group-Ring-Determinant-Conjecture-September-23-2026/build/main.tex\#L107}{thm:input}.
\item OpenAI, \emph{A Torsion-Free Group Algebra That Is Not Directly Finite}, 2026-10-04; citation key \texttt{OAI:A-Torsion-Free-Group-Algebra-That-Is-Not-Directly-Finite-October-4-2026}. \href{https://github.com/openai/math/blob/adc7f1241b42e322a6451854ab7e4b4c146bf78a/preprints/A-Torsion-Free-Group-Algebra-That-Is-Not-Directly-Finite-October-4-2026/direct-finiteness.pdf}{PDF (GitHub)}; \href{https://github.com/openai/math/raw/adc7f1241b42e322a6451854ab7e4b4c146bf78a/preprints/A-Torsion-Free-Group-Algebra-That-Is-Not-Directly-Finite-October-4-2026/direct-finiteness.pdf}{PDF (direct)}; \href{https://github.com/openai/math/blob/adc7f1241b42e322a6451854ab7e4b4c146bf78a/preprints/A-Torsion-Free-Group-Algebra-That-Is-Not-Directly-Finite-October-4-2026/README.md}{Manuscript page and citation}; \href{https://github.com/openai/math/blob/main/preprints/A-Torsion-Free-Group-Algebra-That-Is-Not-Directly-Finite-October-4-2026/direct-finiteness.pdf}{Latest version}; Statement source: \href{https://github.com/openai/math/blob/adc7f1241b42e322a6451854ab7e4b4c146bf78a/preprints/A-Torsion-Free-Group-Algebra-That-Is-Not-Directly-Finite-October-4-2026/build/sections/introduction.tex\#L13}{thm:main}.
\item Family 197: \href{https://github.com/openai/math/blob/adc7f1241b42e322a6451854ab7e4b4c146bf78a/lean/docs/197.md}{Lean scope}.
\item \href{https://github.com/openai/math/blob/adc7f1241b42e322a6451854ab7e4b4c146bf78a/lean/ComparatorChallenges/GroupRingDeterminant.json}{Comparator GroupRingDeterminant.json}.
\item \href{https://github.com/openai/math/blob/adc7f1241b42e322a6451854ab7e4b4c146bf78a/lean/ComparatorChallenges/GroupRingDeterminant.lean}{Comparator GroupRingDeterminant.lean}.
\item \href{https://github.com/openai/math/blob/adc7f1241b42e322a6451854ab7e4b4c146bf78a/lean/ComparatorChallenges/KaplanskyDirectFiniteness.json}{Comparator KaplanskyDirectFiniteness.json}.
\item \href{https://github.com/openai/math/blob/adc7f1241b42e322a6451854ab7e4b4c146bf78a/lean/ComparatorChallenges/KaplanskyDirectFiniteness.lean}{Comparator KaplanskyDirectFiniteness.lean}.
\item \href{https://github.com/openai/math/blob/adc7f1241b42e322a6451854ab7e4b4c146bf78a/lean/ComparatorChallenges/KaplanskyFinitelyPresented.json}{Comparator KaplanskyFinitelyPresented.json}.
\item \href{https://github.com/openai/math/blob/adc7f1241b42e322a6451854ab7e4b4c146bf78a/lean/ComparatorChallenges/KaplanskyFinitelyPresented.lean}{Comparator KaplanskyFinitelyPresented.lean}.
\item \href{https://github.com/openai/math/blob/adc7f1241b42e322a6451854ab7e4b4c146bf78a/lean/ComparatorChallenges/OddKaplansky.json}{Comparator OddKaplansky.json}.
\item \href{https://github.com/openai/math/blob/adc7f1241b42e322a6451854ab7e4b4c146bf78a/lean/ComparatorChallenges/OddKaplansky.lean}{Comparator OddKaplansky.lean}.
\item Family 197: \href{https://github.com/openai/math/blob/adc7f1241b42e322a6451854ab7e4b4c146bf78a/reasoning\_traces/kaplansky-direct-finiteness-characteristic-two.pdf}{Reasoning summary}.
\item \href{https://github.com/openai/math/blob/adc7f1241b42e322a6451854ab7e4b4c146bf78a/CONTENTS.md}{OpenAI catalogue entry}.
\end{itemize}
\end{document}
